\section{The comparison development and finite geodesic runout}
\label{sec:comparison}

The preceding estimates describe two exterior evolutions and a
double-null interior evolution. We now place them in a single vacuum
development containing the complete initial bridge. We then prove two
geometric facts needed for the perturbation argument: a geodesic of
finite future proper duration first leaves this comparison development
through an interior branch or corner, and an open family containing a
regular exiting geodesic contains a branch geodesic. The comparison
development will be an open subspacetime of the full maximal globally
hyperbolic development; we never assert that it exhausts that development.

Throughout the section the Kerr center is fixed, with
$0<\mathfrak a<M$. All depths, finite overlap regions, and constants are
chosen for this center before the final $p_{10}$ smallness threshold.
We first construct a finite Cauchy range around the bridge and use
its overlap collars to attach the late exterior and interior pieces.
We then control geodesic energy through compact and far exterior
excursions, which rules out finite proper-time runout there. Finally,
the interior geodesic equations distinguish branch and corner runout
and let us perturb a regular corner track to a branch.

\subsection{Finite Cauchy ranges and the complete bridge}

We use the signed Killing-time-zero bridge in the two-ended
Boyer--Lindquist extension through the bifurcation sphere. Simultaneous
reversal of the Killing and axial conventions allows either prescribed
future orientation. Write $T_K$ for a continued Killing coordinate and
$T_e=\pm T_K$ for the future-oriented exterior time at an end. Recall
\[
 \Delta=r^2-2Mr+\mathfrak a^2,
 \qquad r_*'=\frac{r^2+\mathfrak a^2}{\Delta},
 \qquad
 \kappa_+=\frac{r_+-r_-}{2(r_+^2+\mathfrak a^2)}>0.
\]
In Cartesian coordinates at either end the reference bridge has
$h_*-\delta=O(r^{-1})$ and connection coefficients $O(r^{-2})$, with
the corresponding differentiated symbol bounds. Consequently, at every
fixed order, its covariant weighted seminorms are equivalent to the
ordinary component symbol seminorms. On compact sets they control the
usual local norms. No asymptotic coefficient of the actual datum is
fixed by this identification.

Fix first the profile depth $q_0$ and common inner rate $\kappa>0$ from
Section~\ref{sec:interior}. The exact initial interior cylinder
$\mathcal S_* =\{Q=q_0\}=\{s=0\}$ is a graph strictly above $r_-$.
Choose the computing boundary $r_b$ strictly below this graph in radial
value, with margins on both sides and $r_b>r_-$. The metric dual of
$ds$ is past timelike; the future side of $\mathcal S_*$ is its
smaller-radius side.

At each end use the horizon-crossing time
\begin{equation}\label{cmp:bent-time}
 t_e=T_e+r_*-G(r).
\end{equation}
Here $G'=1$ through the fixed interior and trapping ranges,
$1\leq G'\leq r_*'$ in an exterior transition, and $G=r_*$ near
infinity. Ordinary ingoing angular coordinates are used at the
corresponding future outer horizon. The temporal sign was established
in the geometric construction; explicitly,
\begin{equation}\label{cmp:temporal-sign}
 \bigl(r^2+\mathfrak a^2\cos^2\vartheta\bigr)
       g_K^{-1}(dt_e,dt_e)
 =-2(r^2+\mathfrak a^2)G'+\Delta(G')^2
       +\mathfrak a^2\sin^2\vartheta<0.
\end{equation}
For $G'=1$ the right side is
$-r^2-2Mr-\mathfrak a^2\cos^2\vartheta$; convexity gives the sign in
the exterior transition.

Moving that transition outward makes $G-r_*$ arbitrarily positive at
each fixed inner radius, since the exterior integral of $r_*'-1$
diverges logarithmically. Moreover $G-r_*\geq0$ outside, and
$(G-r_*)'=1-r_*'>0$ between the horizons where $G'=1$. Thus the two
partial slices $t_e=0$ lie to the future of the bridge, or on it, and
agree with it near their respective infinities. Their portions through
the chosen depth are separated with a fixed margin in $T_K$, including
their subsequent nonnegative-time sides before that depth. Choose the
delay so that the bounded matching bands required at the two ends of
$\mathcal S_*$ are disjoint.

\begin{lemma}[Finite bridge evolution]\label{cmp:finite-range}
Fix the depths and delayed slices just described, together with bounded
forward matching bands on them. After enlarging these fixed regions,
there is a finite vacuum Cauchy development of every sufficiently
$p_{10}$-small datum on the complete bridge which contains the two
delayed slices and the matching bands, as well as a past collar of the
whole bridge. On the delayed slices the induced data are small at
metric order $H^{10}$ and second-form order $H^9$ on compact sets, and
retain the prescribed symbol bounds at the ends. The artificial future
cuts are strict future exits. Higher fixed smooth norms are finite,
without an additional smallness hypothesis.
\end{lemma}

\begin{proof}
We first construct a suitable exact Kerr range. Starting at the bridge,
cut in the future at a radius strictly smaller than $r_b$ but larger
than $r_-$, and at $t_e=L$ on both ends, where the fixed number $L$ is
large enough to leave all required bands and margins. In the trapped
region impose both time inequalities. At either future outer horizon
its own time coordinate is regular and the opposite time inequality is
automatic. Every active cut is a spacelike future exit. We construct on
a slightly larger range and restrict only afterward.

For completeness, the causal control of this exact finite range is as
follows. In the trapped block, $r$ increases strictly to the past. The
Boyer--Lindquist cone inequality gives
\begin{equation}\label{cmp:bl-cone}
 \left|\frac{dT_K}{dr}\right|
 \leq
 \frac{\sqrt{(r^2+\mathfrak a^2)^2
                     -\Delta\mathfrak a^2\sin^2\vartheta}}
      {|\Delta|}.
\end{equation}
This follows by Cauchy--Schwarz on the positive directions tangent to
an $r$ level. Near $r_+$ the right side is $|r_*'|+O(1)$, so the
derivatives of both $r_*+T_K$ and $r_*-T_K$ are bounded above in the
past-directed radial parameter. The positive interior Kruskal variables
\[
 U_+=e^{\kappa_+(r_*+T_K)},\qquad
 V_+=e^{\kappa_+(r_*-T_K)}
\]
therefore have finite nonnegative limits. Indeed each logarithm becomes
monotone after subtracting a bounded multiple of $r$ and hence has a
limit in $[-\infty,\infty)$.

The angular squares in the Kerr metric control
$d\vartheta$ and
$\sin\vartheta((r^2+\mathfrak a^2)d\varphi-\mathfrak a\,dT_K)$.
Replacing $\mathfrak a/(r^2+\mathfrak a^2)$ by its horizon value shows
that the round speed of the corotating angle is
$O(|\Delta|^{-1/2})$ in the radial parameter, an integrable bound.
Together with $U_+,V_+$ this gives a finite point in regular outer
horizon coordinates. To check regularity directly in the
Boyer--Lindquist expression, $r-r_+$ is a signed product $U_+V_+$
times a smooth nonvanishing factor; after corotation, the mixed
$dT_K$--angular coefficients are $O(\Delta)$, and the sum of the
$dT_K^2$ and $dr_*^2$ coefficients is $O(\Delta^2)$. Substitution
therefore gives the smooth nondegenerate horizon metric, up to a smooth
angular-coordinate adjustment.

A past causal continuation at this horizon either follows its generator
toward the bifurcation sphere or enters the corresponding exterior.
There $T_e$ decreases toward its bridge value zero. The exterior radial
causal-speed bound is $O(r-r_+)$ near the horizon and is bounded at
infinity. A curve starting strictly outside cannot reach the horizon or
spatial infinity in its bounded remaining $T_e$ interval. Ordinary
continuation at a finite interior endpoint then carries it to the
bridge. The same argument applies in bounded interior radial ranges.

At bounded spatial radius, the two time cuts give compact closure in
regular coordinates. Compact subsets have compact causal shadows toward
the bridge; far out only bounded-time propagation with bounded speed is
involved. Adding a short open past collar gives an exact Cauchy range
with all the stipulated margins. The boundary signs require only
fixed subextremality.

Apply finite-time vacuum Cauchy stability in wave-map coordinates on an
enlargement of this range. For compact portions include their compact
causal shadows. The $p_{10}$ bound supplies small
$H^{10}\times H^9$ geometric data there. The positive wave-energy flux
across a fixed uniformly spacelike bent surface, after commutation
through order nine, controls tangential metric derivatives through ten
and normal first-jet derivatives through nine. Thus these Cauchy orders
are obtained by a spacelike flux estimate, without a graph-restriction
loss at the highest derivative. Forming the induced metric and second
fundamental form uses exactly those orders. At infinity the bent slices
are the original ends; uniformly local end charts, bounded propagation
speeds, and uniqueness on overlaps give the corresponding finite-time
control there. All higher symbol seminorms remain finite.

The spacelike signs persist by low-order closeness. Past curves cannot
leave through the artificial future cuts. The smooth Cauchy temporal
function on the exact globally hyperbolic development is supplied by
Bernal--S{\'a}nchez \cite[Theorem~1.1]{BernalSanchez2005}. On a fixed
bounded-radius enlargement it remains temporal with a uniform margin
for the perturbed metric. In the far part the exterior times are
temporal, speeds are bounded, and the time to the bridge is bounded.
An excursion from a fixed bounded radius can therefore travel only a
fixed additional distance. Enlarge the compact temporal-function region
beyond that distance and handle curves entirely far out in exterior
time. The resulting coordinate-length bounds preclude an inextendible
endpoint inside these bounded ranges. The short past collar is handled
by the same bounded-geometry construction: every future-inextendible
causal curve there crosses the bridge. These are all fixed finite-time
choices, independent of any later entry or experiment time.
\end{proof}

\begin{proposition}[The comparison development]\label{cmp:development}
There is a single sufficiently small $p_{10}$ neighborhood such that,
for each datum $d$ in it, the preceding exterior and interior solutions
assemble with the finite bridge evolution into a smooth time-oriented
vacuum development $\mathcal P_d$ with the complete bridge $\Sigma$
as a Cauchy surface. It contains an entire spacelike cylinder
$\mathcal S\simeq\mathbb R\times\mathbb S^2$ carrying the entry data
of Section~\ref{sec:interior}, followed by its full product double-null
wedge. There is a time-orientation-preserving open isometric embedding
$\mathcal P_d\hookrightarrow M_d$ into the full MGHD of $d$.
\end{proposition}

\begin{proof}
Evolve the two exterior problems from the geometric data on $t_e=0$.
Choose a fixed matching time $t_a>1$ inside the finite range of
Lemma~\ref{cmp:finite-range}. An enlarged band about $t_e=t_a$ reaches
slightly deeper than every cylinder intersection that will be used.
Both evolutions are vacuum on the overlapping domains.

We describe the identification because its regularity is needed at
$\mathcal S$. Solve a wave map from the exterior evolution into the
finite evolution with the initial slice identification and matching
unit normals. On the fixed earlier slab and matching band it is close
to the identity, with a strict deeper outflow margin and uniform control
at infinity. Pulling the target metric back by this map gives matching
initial first jets: the metric values agree, the second fundamental
form fixes the normal evolution of the spatial components, and the
relative wave-map condition fixes the remaining normal derivatives.
The two metrics then solve the reduced wave equations with a common
target. Ordinary uniqueness on the enlarged domain, with the exiting
boundary signs, proves that the map is an isometry on the required
smaller region.

The fixed-time map has $H^{10}$ wave control on compact sets. After
differentiating its equation, the Christoffel terms need derivatives
through nine. Composing the target coefficients with the map costs
their through-nine spacetime norms on enlarged sets times bounded
composition constants under the $H^{10}$ bootstrap. This follows from
Leibniz and interpolation estimates, bounded low derivatives and inverse
Jacobian, and change of variables. Smooth background compositions cost
the map difference. The resulting spacetime $L^2$ terms have finite-time
$L^1$ bounds; current-order terms are treated by Gronwall. The metric
order-ten estimates therefore give the required smallness for the map
and, after straightening the finite coordinates near the actual matching
face, at least spacetime $H^9$ closeness through the join. Cut off this
small coordinate change farther back. Uniformly local end estimates
give the same low diffeomorphism control near infinity. An embedding on
the smaller region follows from $C^1$ closeness and the fixed margins.
Only a fixed finite interval is involved in this gauge comparison.

In these identified charts join the two prepared end portions of
$\mathcal S=\{s=0\}$ to the central preparation within the bounded
overlaps. Use small smooth interpolations, implemented by small vector
fields for angular maps. The radial positions of the instantaneous
$Q=q_0$ graphs stay uniformly close to the exact graph even at late
times; the instantaneous profile comparison does not introduce an
accumulated angular error in this assertion. In the bounded joining
regions the integrated phases, with their fixed constants, are close
as well.
The timelike conormal and sphere-transversality margins persist.

Retain the finite future piece below $\mathcal S$, with both cuts
$t_e\leq t_a$, and attach each late exterior piece $t_e\geq t_a$
below $\mathcal S$. The deeper computing boundaries are discarded;
all joins lie strictly before them. The agreement bands contain the
necessary cylinder neighborhoods. We identify the same hypersurfaces
in common straightened coordinates, not unequal coordinate-time levels.
This constructs the whole cylinder with the entry estimates already
proved. Across it, geometric Cauchy uniqueness identifies the
double-null solution with matching induced data, using optical
coordinates in local collars.

The resulting topology retains disjoint interiors and identifies common
closed matching faces, with the common collars providing charts.
We do not retain two open future continuations of the same face. The
faces are closed on their respective matching sides and have no extra
finite accumulation points, since late end parameters have unbounded
end times. Each identification is a homeomorphism. At a corner of the
joining prescription, first identify the two lower charts and then
cross $\mathcal S$. These product neighborhoods give a Hausdorff smooth
manifold, together with the original charts and the open past collar.

It remains to check the Cauchy property for the complete bridge. In the
wedge, write
\[
 s=\frac{y_0+y_1}{2},\qquad x=\frac{y_1-y_0}{2}.
\]
Along past causal curves $s$ decreases and $|dx/ds|\leq1$; angular
speeds are bounded on bounded coordinate ranges. Such a curve reaches
$\mathcal S$ or has an ordinary interior continuation before doing so.
In a late exterior part, $t_e$ decreases with finite-time speed bounds.
The sign of $\mathcal S$ excludes a past exit there, so the curve
proceeds into the finite block. Its remaining passage to $\Sigma$ is
controlled by Lemma~\ref{cmp:finite-range}. Curves in the past collar
reach $\Sigma$ to the future, and a crossing of $\Sigma$ has only the
prescribed orientation. Thus every inextendible causal curve meets
$\Sigma$ exactly once. The vacuum MGHD theorem supplies the stated
open embedding, which we use henceforth as an identification.
\end{proof}

The construction has now placed all of our estimates inside $M_d$.
The next objective is to locate the first departure of a finite-duration
geodesic from $\mathcal P_d$. This requires proper-time control, since
bounded spacetime coefficients alone do not bound a geodesic's boost.

\subsection{Compact exterior excursions}

Consider one exterior end, with time $t=t_e$ and co-moving spatial
coordinates $y$. On any fixed compact cylinder containing the matching
interface, the mixed exterior estimates give a majorant
\begin{equation}\label{cmp:compact-integrable}
 \|g-\overline G_{p(t),c(t)}\|_{C^1}
       +|p'(t)|+|c'(t)|\leq\rho_c(t),
 \qquad \rho_c\in L^1([0,\infty)).
\end{equation}
The low errors are uniformly small and the parameters remain small.
The integrability follows by mixed Sobolev embedding and
Cauchy--Schwarz with the static weight $A_S>1$. Constants in
\eqref{cmp:compact-integrable} are independent of the starting time
and duration of an excursion.

Let $U$ be a future unit geodesic tangent. Put
$\dot t=dt(U)>0$, $V=U/\dot t$, and use a prime for differentiation
along the track with respect to $t$. Normalized causal velocities $V$
are bounded on fixed compact cylinders. Define
\[
 T_{\mathrm{co}}=\partial_t|_y,
 \qquad E=-g(T_{\mathrm{co}},U),
 \qquad P=g(\Phi_{p,c},U),
\]
where $\Phi_{p,c}$ is the axial Killing field of the instantaneous
exact Kerr reference. The deformation identity and frozen stationarity
and axisymmetry yield
\begin{equation}\label{cmp:energy-momentum}
 |E'|+|P'|\leq C\rho_c(t)\dot t.
\end{equation}
The time derivative of $\Phi_{p,c}$ is included in the parameter cost.

\begin{lemma}[Compact excursions]\label{cmp:compact-excursions}
For a future unit timelike geodesic which is inextendible in
$\mathcal P_d$ and never crosses $\mathcal S$, every portion
$[t_1,t_2]$ inside a fixed sufficiently
large co-moving radius satisfies
\begin{equation}\label{cmp:compact-speed}
 \dot t(t)\leq C\dot t(t_1),\qquad t_1\leq t\leq t_2,
\end{equation}
with $C$ independent of its duration and of the number of visits to the
outer-horizon strip. If the portion starts and ends at an interface
where $T_{\mathrm{co}}$ is uniformly future timelike, then
\begin{equation}\label{cmp:return-energy}
 E(t_2)\leq E(t_1)
       \exp\left(C\int_{t_1}^{t_2}\rho_c(t)\,dt\right).
\end{equation}
In particular, there is no fixed multiplicative loss at each return.
\end{lemma}

\begin{proof}
We first control the horizon strip. Let $H$ be the future outer-horizon
generator of the base metric $b$, normalized by its Killing parameter,
so $\nabla_HH=\kappa_+H$. Choose a stationary axially invariant future
timelike field $N_0$ there. Invariance gives
\[
 b(\nabla_HN_0,H)=-\kappa_+b(N_0,H)>0.
\]
At the horizon, $dr_K(V)\leq0$ for normalized future causal vectors,
with equality only on the generator; here $r_K$ denotes the reference
Kerr radius. Multiplying $N_0$ by $e^{K_0(r_K-r_+)}$ adds
\[
 2K_0e^{K_0(r_K-r_+)}dr_K(V)b(N_0,V)
\]
to its deformation form. This is positive off the generator. Compactness
of normalized causal directions, followed by continuity, gives a fixed
strip and a stationary future timelike field $N_H$ on it for which
\begin{equation}\label{cmp:redshift}
 (\mathcal L_{N_H}b)(V,V)>c_*>0.
\end{equation}
The strict sign and timelikeness persist for $g$. Hence
$-g(N_H,U)\asymp\dot t$ is nonincreasing during each visit to this
strip.

Below a smaller interior-side level, $dr_K(U)/\dot t<-c$ all the way
to $\mathcal S$. A geodesic continuing while avoiding $\mathcal S$
cannot enter this region: it cannot return through the level and would
reach $\mathcal S$ in bounded $t$. Outside the strip on the exterior
side, stationary energy together with angular momentum controls the
normalized velocity. Indeed in exact Kerr
\[
 b\left(\partial_T+
       \frac{\mathfrak a}{r^2+\mathfrak a^2}\partial_\varphi,
       \partial_T+
       \frac{\mathfrak a}{r^2+\mathfrak a^2}\partial_\varphi\right)
 =-\frac{\Delta(r^2+\mathfrak a^2\cos^2\vartheta)}
          {(r^2+\mathfrak a^2)^2}<0.
\]
On the fixed exterior compact range the coefficients are bounded and
the timelike margin is uniform, also at the axis. Transferring this
combination to each nearby frozen chart and then perturbing gives
\begin{equation}\label{cmp:momentum-control}
 \dot t\leq C(|E|+|P|).
\end{equation}
Thus no timelike-Killing-field assumption is made inside the ergoregion.

Set $X_1=|E(t_1)|+|P(t_1)|$ and
\[
 B(t)=X_1+C\int_{t_1}^t\rho_c(v)\dot t(v)\,dv.
\]
If the portion starts in the strip, include its initial $N_H$ energy
in $X_1$. Outside the strip, \eqref{cmp:energy-momentum} and
\eqref{cmp:momentum-control} give $\dot t\leq CB$. Within the strip,
its nonincreasing $N_H$ energy bounds $\dot t$ by the value at the
initial point or latest strip entry. Using a fixed overlap of the two
control regions and the monotonicity of $B$, the same bound holds
there. Consequently $B'\leq C\rho_c B$, and Gronwall proves
\eqref{cmp:compact-speed}, with one majorant for all strip visits.

At the large interface $X_1\leq CE(t_1)$. Integrating the $E$ equation
itself now gives
\[
 E(t_2)\leq E(t_1)
   \left[1+C_0\left(e^{C\int_{t_1}^{t_2}\rho_c}-1\right)\right]
 \leq E(t_1)e^{C'\int_{t_1}^{t_2}\rho_c}.
\]
This proves the return estimate \eqref{cmp:return-energy}.
\end{proof}

\subsection{Integrability along far portions}

We next bound the geodesic energy between compact excursions. In the
far lab coordinates $(t,z)$ recall the emitter variables
\[
 t=q+r,\qquad z=C_e(q)+rn,\qquad |n|=1,
 \qquad k=1-C_e'(q)\cdot n,
\]
and the field
\[
 T_l=\partial_t|_z,\qquad D=T_l+s_s\cdot\partial_z,
 \qquad g=B+h,\qquad \psi=Dh.
\]
In this subsection $q$ is retarded time and $r$ is emitter radius;
neither is an interior optical variable. The emitter speed is small,
and $k$ is bounded above and away from zero. Interpolate from
$T_{\mathrm{co}}$ near the compact interface through a fixed overlap
to $D$ far out, obtaining a uniformly future timelike field $Z_K$
outside the interface. Its difference from either prescribed field on
the overlap, with first derivatives, has an integrable parameter cost.
Define $E_K=-g(Z_K,U)\asymp\dot t$. The deformation identity gives
\begin{equation}\label{cmp:outer-energy}
 |(\log E_K)'|
       \leq C|(\mathcal L_{Z_K}g)(V,V)|.
\end{equation}

\begin{lemma}[Outer track integrability]\label{cmp:outer-integrability}
Along any one future unit timelike geodesic which is inextendible in
$\mathcal P_d$, eventually proceeds in this end, and never crosses
$\mathcal S$, the right side of
\eqref{cmp:outer-energy} is integrable over all its outer portions,
including arbitrarily many returns to the compact interface.
\end{lemma}

\begin{proof}
We prove separately the integrability of the averaged target terms and
of the field contraction. Besides $\psi(V,V)$ and the already
integrable overlap errors, the deformation cost is bounded by
$C(|DB|+|\partial s_s|)$. Write $x_{\mathrm{av}}$ for the implicit
averaging position, defined by
$x_{\mathrm{av}}=z+\lambda_s(t,x_{\mathrm{av}})$. The target
construction gives $Dx_{\mathrm{av}}=0$, averaging scale
$d(x_{\mathrm{av}})=|x_{\mathrm{av}}|$, and delays in $(0,1/10)$.
Consequently
\begin{equation}\label{cmp:averaged-cost}
 |DB|+|\partial s_s|
 \leq\frac{C}{|x_{\mathrm{av}}|}
       \int_{t-|x_{\mathrm{av}}|/10}^{t}|p'(v)|\,dv.
\end{equation}

Along every late outer portion of this fixed track,
$|x_{\mathrm{av}}|\leq2t$ and $r\leq Ct$, uniformly over repeated
visits. Indeed the far lab speed is at most $1+$ a small constant.
At entry from bounded co-radius one has
$|z|\leq O(1)+\epsilon't$, and past sampling gives
$x_{\mathrm{av}}-z=O(\epsilon't)$ for small drift $\epsilon'$.
Starting from any fixed finite point changes only a bounded initial
time interval. The averaging interval in \eqref{cmp:averaged-cost}
therefore lies at times comparable to $t$. If $\mathcal M$ is the
one-dimensional maximal operator, extending $p'$ by zero before startup
gives the bound
\[
 C\langle t\rangle^{-A/2}
       \mathcal M\bigl(\langle\cdot\rangle^{A/2}|p'|\bigr)(t).
\]
The acceleration belongs to weighted $L^2_A$, and $A=1.36>1$.
The $L^2$ maximal inequality and Cauchy--Schwarz prove integrability.

For the field contraction, first consider $t\geq C_1r$ with $C_1$
large and fixed. Here even the spatial supremum of $\psi$ is
integrable. We include the trace argument to justify evaluation on a
geodesic. On a radial shell $R\leq r\leq2R$, use
$\varrho=\log r$. The strong bulk bounds through two angular rotations,
with the comparisons $W_P\geq cR^{0.26}t^A$ and
$w_{P,j}\geq cR^{-0.03}t^A$, imply
\begin{align}
 \int_{C_1R}^{\infty}t^A
   \|\psi\|_{L^2_{\varrho}H^2_\omega}^2\,dt
       &\leq CR^{-0.26}\mathsf F^2,
       \label{cmp:trace-field}\\
 \int_{C_1R}^{\infty}t^A
   \|\partial_{\varrho}\psi\|_{L^2_{\varrho}H^2_\omega}^2\,dt
       &\leq C(R^{0.03}+R^{-0.26})\mathsf F^2.
       \label{cmp:trace-derivative}
\end{align}
The norms use a slightly enlarged shell, and $C_1$ is increased if
necessary to put all of it in the same timelike region. To see the
derivative count, with $S=r(\partial_t+n\cdot\partial_z)$ one has
$r\partial_{r|t,n}=S-r\partial_q$. Good and full derivatives of
$r\psi$ control the two terms, while the field term accounts for
differentiating $r$. Two rotations and one derivative lie strictly
within the strong commutation budget.

The Hilbert-valued multiplicative one-dimensional trace inequality and
$H^2(\mathbb S^2)\hookrightarrow L^\infty(\mathbb S^2)$ now give
\begin{equation}\label{cmp:trace-supremum}
 \int_{C_1R}^{\infty}t^A
       \sup_{R\leq r\leq2R,\omega}|\psi|^2\,dt
 \leq C\left(R^{-0.26}+R^{-(0.26-0.03)/2}\right)\mathsf F^2
 \leq CR^{-0.115}\mathsf F^2.
\end{equation}
The shell sum converges. Cauchy--Schwarz in time, using $A>1$, yields
\begin{equation}\label{cmp:timelike-cone-integral}
 \int_1^\infty
     \sup_{R_0\leq r\leq t/C_1,\omega}|\psi(t,r,\omega)|\,dt
 <\infty.
\end{equation}
The radial loss $0.03$ still leaves a strictly negative trace power,
which is the margin needed for the shell sum.

On the remaining late outer portions $r\asymp t$. Set $c_3=0.01$.
The point estimates and differentiated gauge identity give
\begin{equation}\label{cmp:far-point}
 \begin{aligned}
 |\psi|&\leq Cr^{-1+e/2}
                   \langle q/r^{\delta_f}\rangle^{-A_H/2},\qquad
 |g-\eta|\leq Cr^{-1+c_3},\\
 |\psi(L,L)|+|\psi(L,e_{\mathrm{scr}})|&\leq Cr^{-2+c_3},
 \end{aligned}
\end{equation}
where $L=(1,n)$, $e_{\mathrm{scr}}$ is a spatial unit screen vector
orthogonal to $n$, and the far exponents satisfy
$\nu=10\delta_f$, $e=0.1\delta_f$. We recall why the improved
projections apply to $\psi$. For $p=T_lh$ and the flat null covector
$\ell_n=(1,-n)$, put
\[
 \mathcal C_\beta(p)
  =\ell_n^\alpha
       \left(p_{\alpha\beta}
                    -\tfrac12\eta_{\alpha\beta}\operatorname{tr}_\eta p\right).
\]
Since $\ell_n^\sharp=-L$, contraction against $L$ or a screen vector
has no trace term and equals respectively $-p_{LL}$ or
$-p_{L\,\mathrm{scr}}$. The far gauge estimate bounds these by good
derivatives, coefficient-suppressed full derivatives, and smaller field
and target terms. Good derivatives cost $O(r^{-2+2\eta_*})$; a product
of $h$ with a full derivative can cost
$O(r^{-2+2\eta_*+e/2})$. The fixed choice $c_3=0.01$ absorbs these
losses, giving the stated $O(r^{-2+c_3})$ bound. Passing from $T_lh$
to $\psi$ uses
\[
 \psi=(1-s_s\cdot n)T_lh+r^{-1}s_s^d b_dZh,
\]
with the bounded angular/radial decomposition coefficients $b_d$ of
the far construction. It adds only good derivatives. Thus no derivative
of an unsmoothed velocity is used in this conversion.

Write $V=(1,\mathbf v)$ on the track. Differentiating the emitter
identity and using causality gives
\begin{equation}\label{cmp:retarded-variation}
 \begin{aligned}
 kq'&=1-\mathbf v\cdot n,\\
 q'&\geq-Cr^{-1+c_3},\\
 |\mathbf v-(\mathbf v\cdot n)n|^2
        &\leq C(|q'|+r^{-1+c_3}),\qquad |q'|\leq C.
 \end{aligned}
\end{equation}
Decomposing $V-L$ into radial and screen parts, the improved
contractions in \eqref{cmp:far-point} imply
\begin{equation}\label{cmp:field-contraction}
 |\psi(V,V)|\leq Cr^{-2+c_3}
              +C|\psi|(|q'|+r^{-1+c_3}).
\end{equation}
Where $|q|\geq r^\nu$, the field itself is integrable, since
\[
 \frac{A_H}{2}(\nu-\delta_f)-\frac e2>0.
\]
It remains to control the $|q'|$ term near the cone.

On a late time dyad $R\leq t\leq2R$, clip $q$ to
$[-C_4R^\nu,C_4R^\nu]$, with $C_4$ containing all values in
$|q|\leq r^\nu$. Extend the clipped function by its upper constant
during compact excursions. Near the interface $r$ is bounded and
$q=t-r\gg C_4R^\nu$, so this extension is continuous. It is uniformly
Lipschitz, even for countably many accumulating excursions: on each
outer component the derivative is bounded, and it matches the same
constant on the complement. Its nonsaturated part has $r\asymp R$.
By \eqref{cmp:retarded-variation} its total negative variation is
$O(R^{c_3})$, while its range has length $O(R^\nu)$. Hence
\begin{equation}\label{cmp:clipped-variation}
 \int_{\substack{R\leq t\leq2R\text{ on outer portions}\\
                        |q|\leq r^\nu}}|q'|\,dt
       \leq C(R^\nu+R^{c_3}).
\end{equation}
Multiplying by $CR^{-1+e/2}$ is summable over dyads, because
$\nu,c_3<1-e/2$. The other two terms in
\eqref{cmp:field-contraction} are integrable directly.
Together with \eqref{cmp:timelike-cone-integral},
\eqref{cmp:averaged-cost}, and the fixed-overlap estimate, this proves
integrability over all outer portions of the track, including all
returns to the compact interface.
\end{proof}

\subsection{Proper time and the first-runout alternatives}

\begin{proposition}[No finite exterior runout]\label{cmp:completeness}
A future unit timelike geodesic from $\Sigma$, inextendible as a
geodesic in $\mathcal P_d$, which never crosses $\mathcal S$ has
infinite future proper duration.
\end{proposition}

\begin{proof}
Such a geodesic must eventually proceed in one end with $t_e\to\infty$.
In the finite piece its far excursions have bounded time and bounded
speed; the remaining bounded-radius range, including the matching
faces, has compact closure in enlarged regular charts. Global
hyperbolicity excludes total causal imprisonment there. There is also
no finite proper-time tangent blowup at an interior position limit:
in a regular temporal chart the causal slopes are bounded, the tail has
finite coordinate length, and the parallel-tangent equation
\[
 \frac{dU^\alpha}{dt}
   =-\Gamma^\alpha_{\beta\gamma}
               \frac{dx^\beta}{dt}U^\gamma
\]
is a linear equation with bounded coefficients. It bounds the tangent
and gives ordinary continuation. The future sign of an end matching
face prevents return to the finite block before crossing
$\mathcal S$. At bounded $t_e$, finite propagation and the same local
argument prevent runout in the late end. This proves $t_e\to\infty$.

Extend $Z_K$ to a smooth future timelike field $X$ throughout this end,
agreeing with $Z_K$ outside and near the compact interface. A partition
of unity gives such an extension by convexity of future timelike cones.
On any fixed finite track interval $[t_1,t_2]$ with outer endpoints,
$F(t)=\log[-g(X,U)]$ is absolutely continuous. Write the open set of
times strictly inside the interface as
$\mathcal O=\bigsqcup_j(a_j,b_j)$. Then
\[
 F(t_2)-F(t_1)
  =\int_{[t_1,t_2]\setminus\mathcal O}F'(t)\,dt
      +\sum_j\bigl(F(b_j)-F(a_j)\bigr).
\]
On this finite interval the series converges absolutely, bounded by
$\int|F'|$. At excursion endpoints $X=Z_K=T_{\mathrm{co}}$,
so \eqref{cmp:return-energy} bounds each increment from above by
$C\int_{a_j}^{b_j}\rho_c$. On the complement,
Lemma~\ref{cmp:outer-integrability} bounds the absolute variation.
The resulting upper bound for $F(t_2)-F(t_1)$ is independent of
$t_2$ and the number of excursions. No deformation bound for the
arbitrary continuation $X$ inside the interface has been used.

Thus $E_K$ remains bounded on all outer portions, and
\eqref{cmp:compact-speed} bounds $\dot t$ on every intervening compact
portion by its entering value. An initial or final portion remaining
entirely compact uses that same supremum estimate. Consequently
$\dot t$ is bounded on the entire late track, and
$\int dt/\dot t=\infty$ as $t\to\infty$.
\end{proof}

In the interior wedge we call the following limiting coordinate
behaviors its two \emph{branches} and its \emph{corner}:
\begin{equation}\label{cmp:runout-types}
 \begin{array}{lll}
 \mathcal B_+:&y_0\longrightarrow y_0^*<\infty,
                     &y_1\longrightarrow+\infty,\\
 \mathcal B_-:&y_1\longrightarrow y_1^*<\infty,
                     &y_0\longrightarrow+\infty,\\
 \mathcal K^+:&y_0\longrightarrow+\infty,
                     &y_1\longrightarrow+\infty.
 \end{array}
\end{equation}
These are runout regimes, not points adjoined to the manifold.

\begin{corollary}[First-runout classification]\label{cmp:classification}
Let a future unit timelike geodesic from $\Sigma$ have finite maximal
future proper duration in $M_d$. Follow its initial connected portion
in $\mathcal P_d$ up to first runout. It crosses $\mathcal S$, and
its first runout has exactly one of the behaviors
\eqref{cmp:runout-types}.
\end{corollary}

\begin{proof}
The first-runout time is no later than its finite maximal time in
$M_d$. Proposition~\ref{cmp:completeness} rules out a track avoiding
$\mathcal S$. In the wedge both $y_i$ increase strictly along future
timelike curves. If both had finite limits, then $s,x$ would be bounded.
Causal slopes in regular temporal charts and compactness of the sphere
would give a position endpoint; the parallel-tangent argument just used
would give ordinary geodesic continuation. Thus at least one limit is
infinite, giving exactly the three listed alternatives.
\end{proof}

First runout can precede the geodesic's maximal time in $M_d$; in that
case its endpoint is an ordinary point of $M_d$. This distinction will
be important in the next argument, which uses an ambient regular chart
but makes no assertion about the whole inner boundary of $M_d$.

\subsection{From a regular corner to a branch seed}

We recall exactly the interior bounds used here. On
$y_0+y_1\geq0$ the metric is
\begin{equation}\label{cmp:wedge-metric}
 \begin{aligned}
 g&=-2a\,dy_0dy_1+
    \gamma_{AB}(d\theta^A-b^A dy_0)(d\theta^B-b^B dy_0),\\
 L_0&=\partial_{y_0}+b^A\partial_{\theta^A},\qquad
 L_1=\partial_{y_1},\qquad
 a=qA,\qquad q=q_0e^{-\kappa(y_0+y_1)}.
 \end{aligned}
\end{equation}
Here $a$ is the positive null factor, not the Kerr rotation parameter;
$q$ now denotes the interior exponential scale. The metric
$\gamma$ and its inverse are uniformly comparable with a fixed round
metric, and $A,A^{-1}$ are uniformly bounded. Set
$\chi_i=\tfrac12\mathcal L_{L_i}\gamma$ on angular tensors and
$\ell_i=L_i\log a$. The angular one-forms $\xi_i$ are characterized by
\[
 \xi_0+\xi_1=d_\theta\log a,
 \qquad \xi_0-\xi_1=-a^{-1}\gamma\,\partial_{y_1}b.
\]
Section~\ref{sec:interior} gives bounded $\xi_i$ and nonnegative
functions $g_i\in L^1(\mathbb R)$ such that
\begin{equation}\label{cmp:corner-low-input}
 |\chi_i|_\gamma+|\ell_i+\kappa|
       \leq C(q+g_i(y_i)),\qquad
 |b|_{\mathrm{round}}\leq C(q+g_0(y_0)).
\end{equation}
All angular statements below are intrinsic to the sphere. We use only
these low bounds and the connection identities
\begin{equation}\label{cmp:null-connections}
 \begin{aligned}
 \nabla^g_{L_i}L_i&=\ell_iL_i,&
 \nabla^g_E L_i&=\chi_i(E)^\sharp+\xi_i(E)L_i,\\
 \nabla^g_{L_i}L_j&=a\xi_i^\sharp
       &&(i\ne j),
 \end{aligned}
\end{equation}
for angular $E$, obtained by Koszul's formula.

The next lemma compares the positive optical variables
$Y_i=e^{-\kappa y_i}$, both of which vanish at a corner.
Near a corner, their speed ratio changes
by a small multiplicative factor that is independent of the initial
unit timelike tangent. A reciprocal change of the two null velocity
components can therefore prevent the variables from vanishing at
the same proper time. This will select a branch while preserving
unit normalization.

\begin{lemma}[Corner ratios]\label{cmp:corner-ratios}
Consider a future unit timelike geodesic starting at a fixed point of
the double-null wedge, and put $Y_i=e^{-\kappa y_i}$. Its maximal
remaining wedge track has finite proper duration, has an angular limit,
and ends in a branch or corner. Each $|\dot Y_i|$ is comparable above
and below to its positive initial value, and its screen velocity is
bounded on that fixed track. Moreover there is
$\delta_c(y_0^{\mathrm{in}},y_1^{\mathrm{in}})$, independent of the
starting unit timelike tangent, with $\delta_c\to0$ as both initial
coordinates tend to $+\infty$, for which
\begin{equation}\label{cmp:ratio-control}
 e^{-2\delta_c}
       \frac{|\dot Y_0^{\mathrm{in}}|}{|\dot Y_1^{\mathrm{in}}|}
 \leq\frac{|\dot Y_0|}{|\dot Y_1|}
 \leq e^{2\delta_c}
       \frac{|\dot Y_0^{\mathrm{in}}|}{|\dot Y_1^{\mathrm{in}}|}.
\end{equation}
If the track ends at the corner, then necessarily
\begin{equation}\label{cmp:corner-necessary}
 e^{-2\delta_c}
       \frac{|\dot Y_0^{\mathrm{in}}|}{|\dot Y_1^{\mathrm{in}}|}
 \leq\frac{Y_0^{\mathrm{in}}}{Y_1^{\mathrm{in}}}
 \leq e^{2\delta_c}
       \frac{|\dot Y_0^{\mathrm{in}}|}{|\dot Y_1^{\mathrm{in}}|}.
\end{equation}
For a general fixed starting point the same conclusions hold with a
finite, not necessarily small, logarithmic-variation constant.
\end{lemma}

\begin{proof}
Dots denote proper-time derivatives. Decompose
\[
 U=\dot y_0L_0+\dot y_1L_1+\Xi,
 \qquad \Xi=\dot\theta-b\dot y_0.
\]
Future timelikeness gives $\dot y_i>0$. For $j=1-i$ define the null
momentum $p_i=-g(U,L_i)=a\dot y_j>0$. Unit normalization and
\eqref{cmp:null-connections} give
\begin{align}
 1+|\Xi|_\gamma^2&=2a\dot y_0\dot y_1,
       \label{cmp:unit-normalization}\\
 \frac{\dot p_i}{p_i}
    &=\ell_i\dot y_i+(\xi_i-\xi_j)(\Xi)
                       -\frac{\chi_i(\Xi,\Xi)}{p_i}.
       \label{cmp:momentum-equation}
\end{align}
Since $|\Xi|_\gamma^2/p_i\leq2\dot y_i$, we obtain
\begin{equation}\label{cmp:log-momentum}
 \begin{aligned}
 \left|\frac{d}{d\tau}\log(e^{\kappa y_i}p_i)\right|
   &\leq C(q+g_i(y_i))\dot y_i+C|\Xi|_\gamma,\\
 e^{\kappa y_i}p_i&=\frac{q_0}{\kappa}A|\dot Y_j|.
 \end{aligned}
\end{equation}
Both $Y_i$ decrease, and \eqref{cmp:unit-normalization} becomes
\[
 1+|\Xi|_\gamma^2
       =\frac{2q_0A}{\kappa^2}|\dot Y_0|\,|\dot Y_1|.
\]
On any initial portion of the remaining track, Cauchy--Schwarz and
monotonicity imply
\begin{equation}\label{cmp:angular-length}
 \int|\Xi|_\gamma\,d\tau
   \leq C\left(\int|\dot Y_0|\,d\tau
                    \int|\dot Y_1|\,d\tau\right)^{1/2}
   \leq C\sqrt{Y_0^{\mathrm{in}}Y_1^{\mathrm{in}}}.
\end{equation}
For each $i$ there are also the bounds
\[
 \int q\dot y_i\,d\tau
       \leq\frac{q_0}{\kappa}Y_0^{\mathrm{in}}Y_1^{\mathrm{in}},
 \qquad
 \int g_i(y_i)\dot y_i\,d\tau
       \leq\int_{y_i^{\mathrm{in}}}^{\infty}g_i(v)\,dv.
\]
Their sum, together with \eqref{cmp:angular-length} and a fixed
constant factor, is a choice of $\delta_c$ in
\eqref{cmp:log-momentum}. It is independent of the initial tangent,
finite at any fixed starting point, and small when both initial
coordinates are large.

The logarithmic variation of each $A|\dot Y_i|$ is at most
$\delta_c$. Taking their ratio cancels $A$, giving
\eqref{cmp:ratio-control}; the bounds for $A^{\pm1}$ give individual
speed comparability. The normalization then bounds $\Xi$ on the
fixed track. Also \eqref{cmp:corner-low-input} implies
$\int |b|_{\mathrm{round}}\dot y_0\,d\tau<\infty$.
Since $\dot\theta=\Xi+b\dot y_0$, the angular track has finite round
length and an angular limit.

An individual $|\dot Y_i|$ has a positive lower bound, so the remaining
proper duration is finite. Both $Y_i$ have nonnegative limits. If both
were positive, the optical coordinates would be bounded; the speed
bounds and angular limit would give ordinary continuation in a regular
chart, contradicting maximality of the wedge track. Thus one or both
limits vanish. In the corner case integration of
\eqref{cmp:ratio-control}, using
$\int |\dot Y_i|\,d\tau=Y_i^{\mathrm{in}}$, gives
\eqref{cmp:corner-necessary}. This proves the assertions, including
the finite-variation version starting at a general fixed point.
\end{proof}

\begin{proposition}[A branch seed in the original open family]
\label{cmp:branch-seed}
Let $d$ belong to the fixed base neighborhood and let
$O\subset T\Sigma$ be open. Suppose $w\in O$ has finite future
proper duration in $M_d$ and its endpoint is reached on a regular
geodesic arc in a $C^2$ extension. Then some seed $w'\in O$ first
runs out of $\mathcal P_d$ at one of its noncorner branches.
The new seed need not have a regular exit in that extension.
\end{proposition}

\begin{proof}
Corollary~\ref{cmp:classification} applies to the original geodesic.
If its first runout is a branch, take $w'=w$. Suppose it is a corner.
At a late point change its tangent to
\begin{equation}\label{cmp:reciprocal-change}
 U_\sigma=(1+\sigma)\dot y_0L_0
           +(1+\sigma)^{-1}\dot y_1L_1+\Xi,
 \qquad \sigma>0.
\end{equation}
The product of its null coefficients and its screen velocity are
unchanged, so it is future unit timelike. Its initial $Y$-speed ratio
is multiplied by $(1+\sigma)^2$. The original corner satisfies
\eqref{cmp:corner-necessary}; a corner for the changed track would
satisfy the same necessary condition with this new ratio. The two
allowable intervals are disjoint once
\begin{equation}\label{cmp:corner-exclusion}
 4\delta_c<2\log(1+\sigma).
\end{equation}
For each fixed $\sigma>0$, a sufficiently late change therefore
produces a branch by Lemma~\ref{cmp:corner-ratios}.

It remains to make one choice of $\sigma$ which returns the changed
geodesic to $O$, uniformly as the change point approaches runout. This
is a statement in regular ambient coordinates, not in the degenerating
optical coordinates. On the fixed original tail $\Xi$ is bounded.
Writing $c=1+\sigma$, direct calculation gives
\begin{equation}\label{cmp:rapidity}
 -g(U,U_\sigma)
 =1+(1+|\Xi|_\gamma^2)
                    \left(\frac{c+c^{-1}}2-1\right).
\end{equation}
The relative rapidity is therefore $O(\sigma)$ uniformly late.

Let $\tau_P$ be first runout from $\mathcal P_d$, and let
$\tau_M=T_d(w)$. If $\tau_P<\tau_M$, the runout state is an
ordinary smooth state in $M_d$. If $\tau_P=\tau_M$, the hypothesis
provides a regular geodesic state in the $C^2$ extension. In either
case choose one precompact product temporal chart $(t,x)$ around that
ambient state, with a convex spatial coordinate ball containing a tail
of the original arc. The original unit tangent is bounded there.
Shrink the chart so that one fixed convex ball of spatial slopes around
the reference slopes consists of future timelike slopes at every point
of the chart.

Choose fixed times $t_0<t_1$ before the runout time such that the
original segment $[t_0,t_1]$ has a compact tube contained in
$\mathcal P_d$. At change points $z_n$ after $t_1$ tending to runout,
\eqref{cmp:rapidity} and the bounded original tangent give an
$O(\sigma)$ change of ambient coordinate phase state. The geodesic
vector field is $C^1$ even for a $C^2$ ambient metric. On a compact
phase-space tube its ODE constants are uniform. Backward integration to
$t=t_0$ therefore exists and changes the state by $O(\sigma)$, with
one constant independent of $n$. Its endpoint $p_n'$ lies in the
fixed interior tube for sufficiently small $\sigma$.

We must also show that this entire backward arc lies in
$\mathcal P_d$. Write $x_0(t)$ for the original spatial path and
choose $\chi\in C^\infty([t_0,t_1])$ equal to one at $t_0$ and zero
near $t_1$. Form a reference path on the early interval by
\[
 x_0(t)+\chi(t)\bigl(x(p_n')-x_0(t_0)\bigr),
\]
and then follow $x_0(t)$ to $z_n$. The early part lies in the fixed
compact tube, the rest in the original track, and the whole path is
timelike for small $\sigma$. It has exactly the endpoints of the
backward-integrated geodesic. Both paths have slopes in the common
convex timelike ball. Spatial linear interpolation at the same chart
time is thus a timelike fixed-endpoint homotopy within the ambient box.

Apply the intrinsic-diamond argument of \Cref{cat:homotopy} to the
globally hyperbolic open subspacetime $\mathcal P_d$. For clarity, the
parameters for which the full path lies inside are open, since the path
is compact. They are closed because such paths lie in the intrinsic
compact causal diamond of their two fixed endpoints. Hence every path
of the homotopy, including the changed arc, lies in $\mathcal P_d$.
No causal convexity in the ambient extension is required.

The changed state at $t=t_0$ is uniformly close to the original one.
Ordinary geodesic flow through the fixed earlier compact segment,
including its transverse crossing of $\Sigma$, then gives a seed in
$O$. The choices have the following order: fix the ambient chart, the
two early slices, their compact tube, and an earlier-flow neighborhood
mapping into $O$; choose one sufficiently small $\sigma>0$ for these
tolerances; finally choose $z_n$ so late that
\eqref{cmp:corner-exclusion} holds. The last choice forces a branch
without changing any preceding tolerance.
\end{proof}

In particular, if $d\in\mathcal B(O,N)$ is a datum in a regular-exit
piece of the closed-test reduction, the preceding proposition provides
a branch seed in the same $O$. Although regular exit need not persist
for that seed, every duration and curvature bound of $A(O,N)$ still
holds along it. The following finite-signal construction will use only
this branch and those closed-test bounds.
